% META: {"id": "TCOMPL-INEG-ME-001", "chapitre": "TCOMPL-INEGALITES", "type_objet": "methode", "capacites_codes": ["C1"], "methodes": ["M1"], "mode_creation": "ex_nihilo", "sources_inspiration": [], "status": "needs_review"}
\begin{fichemethode}{M1}{Construire et interpréter une courbe de Lorenz}

\quandUtiliser{%
  Utiliser cette méthode lorsque l'énoncé fournit une repartition (revenus,
  patrimoines, ventes...) par quantiles et demande de construire ou de lire la
  courbe de Lorenz.
}

\pasApas{%
  \begin{enumerate}
    \item Classer les parts par ordre CROISSANT (des plus modestes aux plus
      eleves).
    \item Calculer les frequences cumulees croissantes de la population ($x$) et
      de la grandeur repartie ($y$) — les deux entre 0 et 1 (ou en \%).
    \item Placer les points $(x_i\,;\,y_i)$, en ajoutant toujours $(0\,;\,0)$
      et $(1\,;\,1)$, et les relier.
    \item Tracer la première bissectrice $y = x$ (droite d'égalité parfaite).
    \item Interpréter un point : « les $100x\,\%$ les plus modestes detiennent
      $100y\,\%$ du total » ; plus la courbe s'ecarte de la bissectrice, plus la
      repartition est inegalitaire.
  \end{enumerate}
}

\exempleRedige{%
  Dans un pays, les 50\,\% les plus modestes detiennent 25\,\% du revenu total,
  et les 80\,\% les plus modestes en detiennent 60\,\%. Placer les points cles
  et interpréter.

  \medskip
  \commentaireMarge{Chaque point se lit « cumule contre cumule ».}%
  La courbe de Lorenz passe par $(0\,;\,0)$, $(0{,}5\,;\,0{,}25)$,
  $(0{,}8\,;\,0{,}6)$ et $(1\,;\,1)$.

  \medskip
  Lecture : le point $(0{,}5\,;\,0{,}25)$ signifie que la moitie la plus modeste
  de la population ne detient que le quart du revenu total — la courbe est
  nettement sous la bissectrice, la repartition est inegalitaire.
}

\pieges{%
  \begin{itemize}
    \item \textbf{Piège 1.} Oublier de trier par ordre croissant avant de
      cumuler : la courbe de Lorenz serait fausse (elle doit rester sous la
      bissectrice).
    \item \textbf{Piège 2.} Cumuler la population mais pas la grandeur (ou
      l'inverse) : les DEUX axes portent des cumuls.
    \item \textbf{Piège 3.} Omettre les points $(0\,;\,0)$ et $(1\,;\,1)$.
  \end{itemize}
}

\verifier{%
  Contrôler que la ligne brisee est croissante, sous la bissectrice, et qu'elle
  relie bien $(0\,;\,0)$ a $(1\,;\,1)$.
}

\sEntrainer{\refExos{M1}}

\end{fichemethode}
